% Standalone problem document, generated from the adjacent README.md.
% Compile with XeLaTeX. Mathematical content is maintained in Markdown.
\documentclass[11pt,a4paper]{article}
\usepackage[fontsize=10.5pt]{scrextend}
\usepackage[margin=22mm,headheight=15pt,headsep=9mm,footskip=11mm]{geometry}
\usepackage{fontspec}
\setmainfont{texgyrepagella}[Extension=.otf,UprightFont=*-regular,BoldFont=*-bold,ItalicFont=*-italic,BoldItalicFont=*-bolditalic]
\setsansfont{texgyreheros}[Extension=.otf,UprightFont=*-regular,BoldFont=*-bold,ItalicFont=*-italic,BoldItalicFont=*-bolditalic]
\setmonofont{texgyrecursor}[Extension=.otf,UprightFont=*-regular,BoldFont=*-bold,ItalicFont=*-italic,BoldItalicFont=*-bolditalic,Scale=MatchLowercase]
\usepackage{amsmath,amssymb}
\usepackage{unicode-math}
\setmathfont{texgyrepagella-math.otf}
\usepackage{xcolor,microtype,enumitem,needspace,fancyhdr}
\usepackage{bookmark}
\definecolor{ink}{HTML}{17344B}
\definecolor{accent}{HTML}{197D80}
\definecolor{muted}{HTML}{596773}
\hypersetup{colorlinks=true,linkcolor=accent,urlcolor=accent,pdftitle={RA-15 - Query
complexity from Schatten to spectral
norms},pdfauthor={Open Problems in Numerical Linear Algebra}}
\urlstyle{same}
\setlength{\parindent}{0pt}
\setlength{\parskip}{5pt plus 1pt minus 1pt}
\linespread{1.04}
\setlength{\emergencystretch}{3em}
\widowpenalty=10000
\clubpenalty=10000
\displaywidowpenalty=10000
\allowdisplaybreaks[1]
\setlist{nosep,leftmargin=1.5em}
\providecommand{\tightlist}{\setlength{\itemsep}{0pt}\setlength{\parskip}{0pt}}
\providecommand{\pandocbounded}[1]{#1}
\pagestyle{fancy}
\fancyhf{}
\fancyhead[L]{\sffamily\footnotesize\color{muted}OPEN PROBLEMS IN NUMERICAL LINEAR ALGEBRA}
\fancyhead[R]{\sffamily\bfseries\footnotesize\color{accent}RA-15}
\fancyfoot[L]{\parbox[b]{0.9\textwidth}{\sffamily\fontsize{8}{10}\selectfont\color{muted}Editorial ratings; a bounded search cannot certify that a problem remains unsolved.\\Literature check: 2026-09-10}}
\fancyfoot[R]{\sffamily\footnotesize\color{muted}\thepage}
\renewcommand{\headrulewidth}{0.3pt}
\renewcommand{\headrule}{\hbox to\headwidth{\color{accent}\leaders\hrule height \headrulewidth\hfill}}
\makeatletter
\renewcommand{\section}{\Needspace{5\baselineskip}\@startsection{section}{1}{0pt}{12pt plus 2pt minus 2pt}{4pt}{\sffamily\bfseries\large\color{ink}}}
\renewcommand{\subsection}{\Needspace{4\baselineskip}\@startsection{subsection}{2}{0pt}{9pt plus 1pt minus 1pt}{3pt}{\sffamily\bfseries\normalsize\color{ink}}}
\makeatother
\setcounter{secnumdepth}{0}
\begin{document}
{\sffamily\small\color{muted}Randomized and low-rank approximation\par}
\vspace{3pt}
{\raggedright\hyphenpenalty=10000\exhyphenpenalty=10000\sffamily\bfseries\fontsize{21}{24}\selectfont\color{ink}Query
complexity from Schatten to spectral norms\par}
\vspace{7pt}
{\sffamily\small\textbf{Difficulty:} extreme\quad\textbf{Importance:} interesting
to the community\par}
{\sffamily\small\bfseries\color{accent}Status: Partially resolved\par}
\vspace{4pt}
{\color{accent}\hrule height 0.8pt}
\vspace{6pt}
\textbf{Rating rationale:} Extreme because the information-theoretic
transition must be uniform in norm parameter, accuracy, and dimension;
community impact is the cost of low-rank approximation across error
measures.

\subsection{Problem statement}\label{problem-statement}

For \(1\le p<\infty\), define the Schatten norm of a real matrix \(M\)
by

\[
\|M\|_{\mathcal S_p}
=\left(\sum_i\sigma_i(M)^p\right)^{1/p},
\]

where \(\sigma_i(M)\) are its singular values; set
\(\|M\|_{\mathcal S_\infty}=\max_i\sigma_i(M)\).

Let \(n\ge2\), \(p\in[1,\infty]\), and \(0<\varepsilon<1/2\). A
randomized algorithm accesses an arbitrary unknown
\(A\in\mathbb R^{n\times n}\) only through exact products \(Ax\) or
\(A^Tx\). Each product costs one query. Query vectors may depend on
previous answers and internal randomness; computation between queries is
unrestricted.

Define \(q_p(n,\varepsilon)\) to be the least integer \(q\) such that
some algorithm using at most \(q\) queries outputs a real unit vector
\(v\) with

\[
\Pr\!\left[
\|A(I-vv^T)\|_{\mathcal S_p}
\le(1+\varepsilon)
\min_{\|u\|_2=1}\|A(I-uu^T)\|_{\mathcal S_p}
\right]\ge\frac{99}{100}
\]

for every \(A\).

\textbf{Open problem.} Determine \(q_p(n,\varepsilon)\) up to universal
constant factors simultaneously in \(p\), \(n\), and \(\varepsilon\),
including the transition when finite \(p\) grows with \(n\) or
\(\varepsilon^{-1}\) and the endpoint \(p=\infty\). Finite-dimensional
saturation must be included, since \(n\) column queries reveal \(A\)
exactly.\\
This square-matrix formulation editorially fixes a two-sided oracle: the
source introduction describes \(Av\), but Algorithm 7.4 uses \(A\) and
\(A^T\). These models coincide on its symmetric lower-bound inputs. The
relative error is in the norm, not its \(p\)th power.

\subsection{Why it matters in numerical linear
algebra}\label{why-it-matters-in-numerical-linear-algebra}

This asks how the matrix-product cost changes between aggregate
singular-value error and worst-direction error.

\newpage
\subsection{References}\label{references}

\begin{enumerate}
\def\labelenumi{\arabic{enumi}.}
\tightlist
\item
  Ainesh Bakshi and Shyam Narayanan,
  \href{https://arxiv.org/html/2304.03191v1\#S1.SS2}{\emph{Krylov
  Methods are (nearly) Optimal for Low-Rank Approximation}},
  arXiv:2304.03191v1 (2023), Open Question 1.11; Theorems 1.1 and 1.5,
  Algorithm 7.4.
\item
  Praneeth Kacham and David P. Woodruff,
  \href{https://drops.dagstuhl.de/entities/document/10.4230/LIPIcs.APPROX/RANDOM.2024.55}{\emph{Faster
  Algorithms for Schatten-\(p\) Low Rank Approximation}}, APPROX/RANDOM
  2024, article 55: later running-time bounds.
\end{enumerate}

\subsection{Status check}\label{status-check}

For fixed finite \(p\), the source gives a dimension-independent query
bound; the spectral lower bound grows with \(\log n\) in its stated
dimension regime. Neither gives optimal uniform dependence on growing
\(p\). The 2024 work improves running time.

Checked 2026-09-08: both arXiv records list only v1. Searches for ``Open
Question 1.11,'' Schatten phase transition, query complexity, and
2025/2026 found no resolution. The older rank-one spectral lower-bound
question was solved by the 2023 source. This is a bounded check.

\subsection{Audit --- 2026-09-10}\label{audit-2026-09-10}

Rechecked \href{https://arxiv.org/html/2304.03191v1}{Bakshi--Narayanan,
Theorem 1.1 and Open Question 1.11}. The spectral endpoint has matching
bounds in its stated dimension regime; optimal growing-\(p\) dependence
remains open. Later Schatten-query searches found no full
characterization; finite-\(p\) near-optimal results and running-time
improvements leave this transition unresolved.
\end{document}
